\documentclass[11pt]{article}

\usepackage[margin=1in]{geometry}
\usepackage{amsmath,amssymb,amsthm,mathtools}
\usepackage{booktabs,longtable,tabularx,array}
\usepackage{enumitem}
\usepackage{microtype}
\usepackage{xcolor}
\usepackage{fancyhdr}
\usepackage[colorlinks=true,allcolors=blue!45!black]{hyperref}

\newtheorem{theorem}{Theorem}
\newtheorem{proposition}[theorem]{Proposition}
\newtheorem{lemma}[theorem]{Lemma}
\newtheorem{corollary}[theorem]{Corollary}
\newtheorem{assumption}[theorem]{Assumption}
\theoremstyle{definition}
\newtheorem{definition}[theorem]{Definition}
\theoremstyle{remark}
\newtheorem{remark}[theorem]{Remark}

\newcommand{\E}{\mathbb E}
\newcommand{\Prb}{\mathbb P}
\newcommand{\cE}{\mathcal E}
\newcommand{\cM}{\mathcal M}
\newcommand{\e}{\mathrm e}
\newcommand{\com}{\mathrm{com}}
\newcommand{\ex}{\mathrm{ex}}
\newcommand{\one}{\mathbf 1}
\newcommand{\pos}[1]{\left[#1\right]_+}

\pagestyle{fancy}
\fancyhf{}
\fancyhead[L]{Order-Level Rescue after Universal Rejection}
\fancyhead[R]{Theory Note}
\fancyfoot[C]{\thepage}
\setlength{\headheight}{14pt}

\title{\bfseries Pay More or Search Wider after Universal Rejection?\\
\large A Complete Theory Note for a Spatial Rescue Mechanism}
\author{}
\date{28 August 2026}

\begin{document}
\maketitle

\begin{abstract}
This note develops and audits an order-level rescue mechanism for a
representative ride request.  The platform commits to an anonymous
transaction-price path $(p_1,p_2)$ and, when allowed, a rescue-footprint
multiplier $s$.  First-window incumbents have already seen the focal order and
may reject strategically in anticipation of the same order's rescue terms.
Universal rejection therefore screens incumbent costs, becomes public
information for the rider, and triggers competition among retained rejectors,
time-homogeneous fresh arrivals, and newly reached outer drivers.  Every
policy is evaluated at its induced cutoff weak perfect Bayesian equilibrium
(WPBE) before the platform optimizes.

Under the maintained uniform-Poisson, type-independent-retention, anonymous-
lottery model, each feasible policy has a unique incumbent cutoff and a unique
completion outcome, up to payoff-equivalent action labels.  We derive an exact
rider aggregation, prove that the incremental values of fixed-footprint
replenishment and expanded reach vanish in both the thin and thick limits, and
show that any strictly positive gap must attain an interior maximum.  A closed-
form ordering of the two peaks is valid in a frozen-hazard Poisson benchmark,
but it is not a general theorem of the full WPBE model; a maintained-primitives
fixed-menu counterexample reverses the order.  Finally, a failure-contingent
committed-reach capacity cost yields a finite global no-search threshold under
a uniform regularity condition and a complete-WPBE local search index.  An
executed-notification cost does not generally have this property at a rider
rescue-activation boundary.

The note also records the exact contribution boundary after an adversarial
literature search through 28 August 2026.  Two-round radius adjustment, joint
pricing and spatial matching, strategic driver rejection, and rejection-aware
notifications each exist in prior work.  The defensible object is their
interaction at a focal order's public universal-rejection node.
\end{abstract}

\section{Research object and audited contribution boundary}

The model studies a constrained posted mechanism, not an unrestricted direct
mechanism.  Before the order is posted, the platform commits to an anonymous
menu that may depend on the public event of universal first-window rejection,
but not on realized driver counts, retention realizations, or private costs.
The central question is when recovery should rely on a higher common second-
window transaction price, a wider notification footprint, or both.

The closest literatures meet at the post-failure node but abstract from
different sides of it.  Hu, Hu, and Zhu~\cite{HuHuZhu2022} provide the closest
temporal-pricing benchmark: they study a two-period price path, strategic
temporal responses, and incoming drivers, but explicitly remove matched-driver
rejection from their main mechanism.  Qin, Yang, and Liu~\cite{QinYangLiu2025}
already propose adaptive two-round broadcasting radii, so expanded radius
after first-round nonresponse is not new.  Ekbatani et al.~\cite{LyftI2026,LyftII2026}
optimize non-exclusive notification sets under rejection uncertainty and
study long-run thickness, while explicitly holding driver acceptance behavior
fixed under the notification intervention.  The new working paper by Wang,
Zhang, Feng, and Ke~\cite{WangZhangFengKe2026} studies pricing and matching
under broadcasting with endogenous distance-sensitive driver response.  It
eliminates any broad claim that joint price and broadcast-footprint design is
new.  Its full text remains the highest-priority manual overlap check.

\begin{longtable}{@{}p{0.20\textwidth}p{0.34\textwidth}p{0.40\textwidth}@{}}
\caption{Adversarial closest-paper map.}\label{tab:lit}\\
\toprule
Paper & Central overlap & Main missing element relative to this note \\
\midrule
\endfirsthead
\toprule
Paper & Central overlap & Main missing element relative to this note \\
\midrule
\endhead
Hu et al. (2022) & Two-period price, strategic temporal response, incoming
drivers & No focal-order rejection screening or rescue footprint \\
Qin et al. (2025) & Two-round broadcast radii and heterogeneous acceptance &
No announced $p_1\to p_2$ same-order holdout equilibrium in the public record \\
Ekbatani et al. (2026) & Rejection-aware notifications, fresh availability,
and long-run thickness & Acceptance response is stochastic/exogenous, not a
price-induced WPBE \\
Wang et al. (2026) & Broadcasting with endogenous distance response; pricing
and matching & Public abstract is one-shot/steady-state, not
failure-contingent same-order rescue \\
Chen and Hu (2020) & Forward-looking two-sided pricing with Poisson arrivals &
Market-level waiting, not rejection of a focal order \\
Castro et al. (2022) & Strategic decline and request reoffer & Priority/reoffer
mechanism, not a common $p_2$ plus spatial scope \\
Xu and Yan (2026) & Joint state-dependent spatial dispatch and pricing &
Binding dispatch/MDP, not nonbinding broadcast and holdout \\
\textbf{This note} & Same-order holdout, fresh supply, rider failure decision,
and reach & Cutoff-WPBE nested inside posted $(p_1,p_2,s)$ design \\
\bottomrule
\end{longtable}

The resulting contribution claim is intentionally joint rather than
ingredient-by-ingredient:

\begin{quote}
We study an order-level recovery mechanism in which public universal
rejection simultaneously screens incumbent private costs, governs rider
continuation, and triggers a common second-window transaction price that pools
retained rejectors with fixed-footprint fresh arrivals and, when chosen, a
newly reached outer annulus.  The platform designs $(p_1,p_2,s)$ subject to the
induced cutoff-WPBE.
\end{quote}

We do \emph{not} claim the first two-period ride-hailing price, the first
strategic driver rejection, the first two-round broadcast, the first matching-
radius decision, or the first joint pricing-and-spatial-matching problem.

\section{Environment}

\subsection{Prices, types, and local thickness}

There is one representative order and two equal-length response windows.
Rider value and drivers' persistent order costs are independent and uniform:
\[
v,c\sim U[0,1].
\]
The rider and incumbent delay factors are $\beta,\delta\in(0,1]$.

The variable $p_j$ has one economic meaning throughout: it is the anonymous
full-pass-through transaction price in window $j$.  The rider pays $p_j$ and
the assigned driver receives $p_j$.
Thus the current two-price model is neither a platform-profit model nor a
separate fare-and-wage model.  Separating the rider fare and driver wage would
enlarge the mechanism to at least four monetary instruments.  The platform's
objective below values completion and charges a search-resource cost.

The first-window incumbent cohort and second-window fresh core cohort are
independent:
\[
N_1\sim\operatorname{Pois}(m),\qquad
N_2^C\sim\operatorname{Pois}(m),\qquad N_1\perp N_2^C.
\]
Here $m$ is a \emph{local per-window driver intensity} faced by a
representative order.  Because simultaneous orders and overlapping
notifications are not modeled, $m$ should not be interpreted as system-wide
market thickness without a separate sparse/mean-field foundation.

\subsection{Supply architectures}

We compare three supply environments.

\begin{enumerate}[label=\textup{\Roman*.},leftmargin=2.4em]
\item \textbf{Retained incumbents only} ($I$): the terminal pool contains
first-window rejectors who remain available.  The platform chooses
$(p_1,p_2)$.
\item \textbf{Fixed footprint with fresh arrivals} ($A$): the same retained
rejectors compete with an independent, time-homogeneous fresh core cohort in
the second window.  The platform chooses $(p_1,p_2)$.
\item \textbf{Expanded rescue search} ($E$): the fixed-footprint fresh cohort
is present under both repeat and rescue, and rescue can additionally reach an
outer annulus.  The platform chooses $(p_1,p_2,s)$.
\end{enumerate}

Architectures $I$ and $A$ are supply counterfactuals, not nested mechanism
sets.  Architecture $A$ is nested in $E$ because $s=1$ exactly reproduces the
fixed footprint.

Let $\omega\in[0,1]$ be the probability that a first-window rejector remains
physically available and eligible for the focal order's terminal lottery.
Retention is independent of cost and of the rider's terminal action.  The
same $\omega$ enters the aggregate pool and the focal incumbent's continuation
payoff.

\subsection{Search multiplier and winner-paid pickup cost}

Normalize the core radius to $R_0$.  The search control $s\in[1,\bar s]$ is an
\emph{area multiplier}:
\[
R(s)=R_0\sqrt{s}.
\]
Spatial Poisson independent increments therefore give an outer fresh cohort
\[
N_2^O(s)\sim\operatorname{Pois}(m(s-1)),
\qquad N_2^O(s)\perp N_2^C.
\]
Let $u=(\ell/R_0)^2$ denote area rank.  A selected outer driver pays the extra
pickup cost
\[
d(u)=\tau(\sqrt u-1)_+,
\qquad \tau\ge0,
\]
after assignment.  Core pickup cost is absorbed into $c$.  This timing is a
maintained assumption: an unassigned notified driver pays no sunk relocation
cost.  If pickup or activation cost were paid before assignment, fresh entry
would require a congestion fixed point rather than the Poisson thinning below.

\subsection{Timing and information}

\begin{enumerate}[leftmargin=2em]
\item The platform publicly commits to the feasible anonymous menu.
\item A rider observes $v$ and posts iff $v\ge p_1$.
\item The $N_1$ incumbents receive the focal order.  Each observes the order,
$(p_1,p_2,s)$, and own $c$, but not other costs or realized driver counts.
\item Incumbents simultaneously accept at $p_1$ or reject and wait.  If at
least one accepts, the order is assigned by an anonymous lottery.
\item Universal rejection is publicly observed.  The rider does not observe
the number of first-window drivers, retention, future arrivals, or volunteers.
She chooses abandon, repeat, or rescue using equilibrium beliefs.
\item Only after the terminal action is selected do fresh drivers first see
the order.  They volunteer when the assigned-order surplus is nonnegative.
Retained incumbents and fresh volunteers enter a common anonymous lottery.
\end{enumerate}

The platform may condition on public universal rejection but not on hidden or
realized counts.  The reach commitment must be verifiable or technologically
enforceable.

\section{Conditional supply and assignment}

For a generic cost distribution $F$, let $a$ be a proposed first-window
incumbent cutoff.  A retained incumbent can volunteer at terminal price $p_j$
only when $a<c\le p_j$, so define
\[
I_j^\omega(a)=\omega m\,[F(p_j)-F(a)]_+.
\]
Under the maintained uniform distribution, write $F(x)=[x]_0^1$.

Define the fresh willing-driver intensities
\[
e_j^I=0,
\qquad
e_j^A=mF(p_j),
\]
and, in the expanded architecture,
\[
e_1^E=mF(p_1),
\qquad
e_2^E=mF(p_2)+m\sigma_{p_2}(s),
\]
where
\[
\sigma_p(s)=\int_1^s F(p-d(u))\,du.
\]
The total terminal willing-driver intensity and coverage are
\[
\Lambda_j^r(a)=I_j^\omega(a)+e_j^r,
\qquad
C_j^r(a)=1-\e^{-\Lambda_j^r(a)}.
\]

\begin{lemma}[Conditional Poisson supply decomposition]
\label{lem:poisson}
Under a first-window cutoff $a$, conditioning on universal rejection leaves
the number of terminal-eligible rejectors Poisson with mean
$I_j^\omega(a)$.  Independent fresh core and outer-annulus volunteers add
their respective Poisson means.  Hence the terminal volunteer count is
Poisson with mean $\Lambda_j^r(a)$ in architecture $r$.
\end{lemma}

\begin{proof}
Mark $N_1$ by driver cost.  Poisson splitting makes the counts with $c\le a$
and $a<c\le p_j$ independent, with means $mF(a)$ and
$m[F(p_j)-F(a)]$.  Universal rejection sets the first count to zero but does
not change the second count's conditional distribution.  Independent
$\omega$-retention thins the latter to mean $I_j^\omega(a)$.  The fresh core
flow is independent.  Spatial independent increments and cost thinning give
outer mean $m\sigma_{p_2}(s)$.  Sums of independent Poisson variables are
Poisson with the sum of their means.
\end{proof}

If a focal volunteer faces $N\sim\operatorname{Pois}(z)$ opponents in an
anonymous lottery, her expected assignment probability is
\[
h(z)=\E\left[\frac1{N+1}\right]
=\frac{1-\e^{-z}}{z},\qquad h(0)=1.
\]
An outer fresh driver's expected payoff is
\[
h(\Lambda_2)[p_2-c-d(u)].
\]
Because the pickup cost is assignment-contingent, $h(\Lambda_2)>0$ scales the
payoff without changing its sign.  Thus volunteering is equivalent to
$c+d(u)\le p_2$, which validates the thinning formula for $\sigma_{p_2}(s)$.

\section{Rider continuation and exact aggregation}

After universal rejection, a rider of value $v$ compares
\[
0,\qquad C_1(a)(\beta v-p_1),\qquad
C_2(a)(\beta v-p_2).
\]
Let $\eta_j(a)$ be the conditional mass of posted riders selecting terminal
action $j\in\{0,1,2\}$.  Ties at a single rider type are immaterial.  If two
terminal actions are payoff-identical for a positive measure, we use the
deterministic convention abandon before repeat before rescue.  This convention
selects labels only; completion and incumbent waiting payoffs coincide across
payoff-identical terminal actions.

Set
\[
P=1-p_1,\qquad q=\e^{-ma},\qquad
\rho_j=1-\frac{p_j}{\beta},\qquad A_j=\rho_j C_j.
\]

\begin{lemma}[Exact rider aggregation]
\label{lem:rider}
Suppose $0\le C_1\le C_2\le1$.  If $C_2>C_1$, the repeat-rescue crossing is
\[
v^{12}=\frac{C_2p_2-C_1p_1}{\beta(C_2-C_1)},
\qquad
v^{12}-\frac{p_2}{\beta}
=\frac{C_1(p_2-p_1)}{\beta(C_2-C_1)}\ge0.
\]
On the active domain $p_2\le\beta$, rescue has positive mass iff $A_2>A_1$
and
\[
P\eta_2=\frac{[A_2-A_1]_+}{C_2-C_1}.
\]
On the full price domain, terminal completion integrated over the original
rider population is
\[
P(\eta_1C_1+\eta_2C_2)=\max\{0,A_1,A_2\}.
\]
Consequently every equilibrium completion outcome satisfies
\begin{equation}
M=P(1-q)+q\max\{0,A_1,A_2\}.
\label{eq:completion}
\end{equation}
If $C_2=C_1$ and $p_2>p_1$, repeat dominates rescue.  If both coverage and
price are equal, action labels are tie-rule dependent but completion is not.
\end{lemma}

\begin{proof}
The three continuation payoffs are affine in $v$ and their upper envelope has
the order abandon, repeat, rescue.  On a strict rescue branch, riders in
$(v^{12},1]$ rescue and direct algebra gives
\[
1-v^{12}=\frac{A_2-A_1}{C_2-C_1}.
\]
Moreover,
\[
C_1\left(v^{12}-\frac{p_1}{\beta}\right)
+C_2(1-v^{12})=A_2.
\]
If rescue is inactive, repeat contributes $[A_1]_+$; if repeat is also
inactive, terminal completion is zero.  Multiplying by the first-window
failure probability yields \eqref{eq:completion}.
\end{proof}

For later reference, expected executed outer notifications on a strict rescue
branch are
\begin{equation}
Q^{\ex}=Pq\eta_2m(s-1)
=qm(s-1)\frac{[A_2-A_1]_+}{C_2-C_1}.
\label{eq:qex}
\end{equation}
The ratio in \eqref{eq:qex} is not defined by continuity at a joint
price-coverage degeneracy; this is one reason not to use executed contacts as
the main cost primitive.

\section{Cutoff-WPBE}

Under a proposed aggregate cutoff $a$, a type-$c$ incumbent's accept-minus-
wait payoff in architecture $r$ is
\begin{equation}
D_r(c;a)=h(ma)(p_1-c)
-\delta\omega\e^{-ma}
\sum_{j=1}^2\eta_j^r(a)h(\Lambda_j^r(a))[p_j-c]_+.
\label{eq:D}
\end{equation}
The first term incorporates first-window assignment competition.  The second
requires universal rejection by other incumbents, survival, rider
continuation, and terminal assignment.

\begin{definition}[Cutoff-WPBE]
A symmetric cutoff $a\in[0,p_1]$, together with rider and fresh-driver
strategies, is a cutoff-WPBE if: (i) incumbents accept for $c<a$ and reject for
$c>a$, with the appropriate weak inequalities at $a$; (ii) riders and fresh
drivers are sequentially rational; and (iii) on-path beliefs after universal
rejection are generated by the cutoff and Lemma~\ref{lem:poisson}.
\end{definition}

\begin{theorem}[Unique incumbent cutoff and completion outcome]
\label{thm:unique}
Fix a feasible policy in $I,A$, or $E$.  Under uniform driver costs,
type-independent retention, cutoff-independent fresh intensities, and the
anonymous lottery, there is exactly one incumbent cutoff
$a^*\in[0,p_1]$.  The cutoff and completion outcome are continuous in policy
and primitives when fresh intensities are continuous.  Full strategies are
unique only up to payoff-equivalent rider labels and measure-zero driver ties.
\end{theorem}

The proof is in Appendix~\ref{app:unique}.  The key inequality is that at every
interior cutoff root, the immediate-acceptance residual crosses zero strictly
downward.  The restrictions in the theorem are substantive.  Type-dependent
retention, pre-assignment entry costs, priority assignment, or general $F$
may destroy the derivative ordering and require a new equilibrium analysis.

\begin{corollary}[Flat price]
\label{cor:flat}
For $p_1=p_2=p$ and $s=1$, the unique incumbent cutoff is $a=p$ in every
architecture.
\end{corollary}

\begin{proof}
For $a<p$, terminal service has the same price and weakly worse timing and
assignment odds than immediate acceptance.  Formally, the residual factors as
$(p-a)$ times a strictly positive bracket because $h(ma)\ge\e^{-ma}$,
$\delta\omega\le1$, and terminal competition cannot improve the focal
driver's assignment probability.  Thus no $a<p$ is a root; at $a=p$, all
lower types accept and higher types have no positive surplus.
\end{proof}

\section{Outer posted-mechanism design}

The feasible policy sets are compact:
\[
\cM_I=\cM_A=\{(p_1,p_2):0\le p_1\le p_2\le1\},
\]
\[
\cM_E=\{(p_1,p_2,s):0\le p_1\le p_2\le1,
\ 1\le s\le\bar s\}.
\]
Policies with active delayed service may impose $p_2\le\beta$ without loss;
the inactive closure is retained for existence.

Free, unbounded reach would make completion weakly favor the largest feasible
scope.  We therefore distinguish the driver's winner-paid pickup cost from a
platform search resource.  The main specification prices the capacity needed
to honor the promised outer footprint after first-window failure:
\begin{equation}
Q^{\com}=(1-p_1)\e^{-ma}m(s-1).
\label{eq:qcom}
\end{equation}
This cost is failure-contingent but does not vanish merely because the rider is
exactly indifferent at the rescue-activation boundary.  It can represent
reserved notification bandwidth, dispatch attention, or the shadow cost of
shared capacity omitted from the representative-order partial equilibrium.
Core notification infrastructure is normalized to zero incremental cost.

For completion value $B>0$ and per-unit capacity cost $\kappa\ge0$, define
\[
J_r^{\com}(\mu,a)=BM_r(\mu,a)-\kappa Q_r^{\com}(\mu,a),
\qquad Q_I^{\com}=Q_A^{\com}=0.
\]
The equilibrium-constrained value is
\begin{equation}
\mathcal V_r(m;\kappa)
=\max_{\mu\in\cM_r}
\min_{a\in\cE_r^{\mathrm{WPBE}}(\mu;m)}J_r^{\com}(\mu,a).
\label{eq:outer}
\end{equation}
Theorem~\ref{thm:unique} makes the inner set a singleton, but writing the
problem this way makes clear that the platform does not optimize at a fixed
cutoff.  Computationally, one may fix $p_1$, optimize $p_2$ and $s$ while
re-solving the WPBE at every evaluation, and then optimize $p_1$.

The objective is completion net of a search-resource cost.  It is not
platform profit or total welfare.  A profit model would separate rider fare
and driver wage and would also include pickup-time cancellation or other
service costs.

\section{What market thickness guarantees}

Let $V_I(m)$ and $V_A(m)$ denote optimized completion in architectures $I$
and $A$.  Define
\[
\Delta_A(m)=B[V_A(m)-V_I(m)],
\qquad
\Delta_E(m;\kappa)=\mathcal V_E(m;\kappa)-BV_A(m).
\]
The first difference compares two supply environments and need not be
nonnegative by set inclusion.  The second is nonnegative because $s=1$
reproduces architecture $A$ at zero outer-capacity cost.

\begin{theorem}[Endpoint collapse under WPBE and full reoptimization]
\label{thm:endpoint}
For finite $\bar s$ and compact policy sets,
\[
\lim_{m\downarrow0}V_I(m)=\lim_{m\downarrow0}V_A(m)=0,
\qquad
\lim_{m\downarrow0}\mathcal V_E(m;\kappa)=0,
\]
and
\[
\lim_{m\to\infty}V_I(m)=\lim_{m\to\infty}V_A(m)=1,
\qquad
\lim_{m\to\infty}\mathcal V_E(m;\kappa)=B.
\]
Consequently both incremental gaps vanish as $m\downarrow0$ and
$m\to\infty$.  For all sufficiently large $m$,
\[
0\le\Delta_E(m;\kappa)
\le B[1-V_A(m)]
\le B\left(\frac{\log m}{m}+\frac1m-
\frac{\log m}{m^2}\right).
\]
If $\kappa>0$, any expanded-search optimizer satisfies
\[
Q_m^{\com,*}
\le\frac{B[1-V_A(m)]}{\kappa}
=O\!\left(\frac{\log m}{\kappa m}\right).
\]
\end{theorem}

\begin{proof}
Completion requires at least one available driver in a reached cohort whose
total mean is no more than $(1+\bar s)m$.  Hence completion is at most
$1-\e^{-(1+\bar s)m}$, and the optimized nonnegative objective vanishes in
the thin limit.

For $m>\e$, choose the feasible flat policy
\[
p_1=p_2=p_m=\frac{\log m}{m},\qquad s=1.
\]
Corollary~\ref{cor:flat} gives $a=p_m$.  First-window completion alone is
\[
\left(1-\frac{\log m}{m}\right)
\left(1-\frac1m\right)\longrightarrow1,
\]
and this policy uses no outer capacity.  Completion is at most one, proving
the thick limits.  Subtraction yields endpoint collapse.  Nesting gives
$\Delta_E\ge0$, and the displayed flat-policy lower bound yields the rate.
Optimality relative to $s=1$ implies
$BM_E-\kappa Q_m^{\com,*}\ge BV_A$; using $M_E\le1$ gives the resource bound.
\end{proof}

\begin{corollary}[Existence, not uniqueness, of an interior peak]
\label{cor:peakexist}
If $\Delta_E(m_0;\kappa)>0$ at some finite $m_0$, then it has a positive
global maximizer in $(0,\infty)$.  The same statement holds for $\Delta_A$ if
$\Delta_A(m_0)>0$.  Neither conclusion implies a unique peak.
\end{corollary}

\begin{proof}
Continuity follows from Theorem~\ref{thm:unique}, the deterministic tie
closure, and Berge's maximum theorem.  Choose thin and thick endpoints at
which the gap is less than half its value at $m_0$.  A maximum on the
intervening compact interval cannot occur at either endpoint.
\end{proof}

This theorem gives a robust scarce-intermediate-thick support result.  It does
not by itself identify the shape inside that support.

\section{The two-peak question}

\subsection{A valid frozen-hazard theorem}

The clean closed forms arise only when normalized hazards and activation
weights are held fixed as $m$ changes.  This is a useful Poisson benchmark,
not an endogenous full-design WPBE comparative static.

\begin{theorem}[Ordered peaks in a frozen-hazard Poisson benchmark]
\label{thm:frozen}
Let $A(m)>0$ be differentiable, nonincreasing, and log-concave, and let
\[
0\le\lambda_I<\lambda_A<\lambda_E
\]
be constants.  Define gross adjacent gains
\[
\Delta_A^0(m)=A(m)\e^{-\lambda_I m}
\left(1-\e^{-(\lambda_A-\lambda_I)m}\right),
\]
\[
\Delta_E^0(m)=A(m)\e^{-\lambda_A m}
\left(1-\e^{-(\lambda_E-\lambda_A)m}\right).
\]
If both gains vanish at the endpoints, each is strictly log-concave with a
unique interior maximizer, and
\[
\arg\max_m\Delta_E^0(m)<\arg\max_m\Delta_A^0(m).
\]
If $A(m)=\e^{-\varphi m}$ and $\varphi+\lambda_I>0$, the maximizers are
\[
m_A^*=\frac{\log[(\varphi+\lambda_A)/(\varphi+\lambda_I)]}
{\lambda_A-\lambda_I},
\]
\[
m_E^*=\frac{\log[(\varphi+\lambda_E)/(\varphi+\lambda_A)]}
{\lambda_E-\lambda_A},
\]
with
\[
m_E^*<\frac1{\varphi+\lambda_A}<m_A^*.
\]
\end{theorem}

\begin{proof}
For $d>0$, $\log(1-\e^{-dm})$ has strictly negative second derivative.
Hence both gains are strictly log-concave.  Put
$d_A=\lambda_A-\lambda_I$ and $d_E=\lambda_E-\lambda_A$.  Their log-slope
difference is
\[
(\log\Delta_A^0)'-(\log\Delta_E^0)'
=\frac{d_A}{1-\e^{-d_A m}}
-\frac{d_E}{\e^{d_E m}-1}>0,
\]
because
\[
\frac{d_A}{1-\e^{-d_A m}}>\frac1m>
\frac{d_E}{\e^{d_E m}-1}.
\]
At the search peak, the search log slope is zero, while the arrival log slope
is positive; the arrival peak lies to its right.  Under exponential
activation, each gain is a difference of two exponentials.  Solving its first-
order condition yields the two closed forms.  The middle inequality follows
from
$\log(b/a)=\int_a^b t^{-1}dt$.
\end{proof}

\subsection{Why the full WPBE does not inherit the order automatically}

In the full model, prices, cutoffs, rider action masses, and per-unit hazards
typically vary with $m$ and differ across supply architectures:
\[
p_j=p_j^r(m),\qquad a=a^r(m),\qquad
\eta_j=\eta_j^r(m),\qquad \lambda_j=\lambda_j^r(m).
\]
Theorem~\ref{thm:frozen} therefore does not prove an ordered peak after
equilibrium response and reoptimization.

\begin{remark}[Maintained-primitives fixed-menu reversal]
The absence of a primitive theorem is not merely abstract.  At
\[
(\beta,\delta,\omega,\tau,\kappa)
=(.1510729,.4580650,.1790836,.8263075,.0100430)
\]
and the fixed menu
\[
(p_1,p_2,s)=(.0371029,.0864844,2.4731078),
\]
re-solving the architecture-specific cutoff-WPBE at every $m$ gives
\[
m_A^*=12.59214,\quad \Delta_A(m_A^*)=.167859,
\]
\[
m_E^*=15.50563,\quad \Delta_E(m_E^*)=.002470.
\]
Thus $m_E^*>m_A^*$, and the search gap changes sign and has multiple lobes.
These gaps use the archived executed-contact objective.  The solver-certified
counterexample lies inside the maintained primitives and is sufficient to make
the corresponding universal fixed-menu peak-order claim unsafe.  It is not an
interval proof, a symbolic characterization of when reversal occurs, or a
counterexample to every alternative search-cost specification.
\end{remark}

One can nevertheless certify an ordering for a particular calibrated model
by a uniform separation condition.

\begin{proposition}[Uniform separation certificate]
Let a compact set $K$ contain the reference peaks, all full-gap maximizers,
and a separator $\bar m\in(m_E^0,m_A^0)$.  Define
\[
\Gamma_A=\Delta_A^0(m_A^0)-
\sup_{m\in K:m\le\bar m}\Delta_A^0(m)>0,
\]
\[
\Gamma_E=\Delta_E^0(m_E^0)-
\sup_{m\in K:m\ge\bar m}\Delta_E^0(m)>0.
\]
If
\[
\|\Delta_A-\Delta_A^0\|_{\infty,K}<\Gamma_A/2,
\qquad
\|\Delta_E-\Delta_E^0\|_{\infty,K}<\Gamma_E/2,
\]
then every full search-gap maximizer is below $\bar m$, and every full
arrival-gap maximizer is above $\bar m$.
\end{proposition}

\begin{proof}
At $m_A^0$, the full arrival gap is at least
$\Delta_A^0(m_A^0)-\varepsilon_A$.  Anywhere weakly below $\bar m$, it is at
most the displayed restricted supremum plus $\varepsilon_A$.  The former is
strictly larger when $2\varepsilon_A<\Gamma_A$.  The search argument is
symmetric.
\end{proof}

This proposition is a robustness certificate, not a primitive sufficient
condition.  It should be verified numerically rather than advertised as the
main source of economic intuition.

\section{Costly search and the no-expansion boundary}

\subsection{Global committed-reach threshold}

Let $\widehat M(p_1,p_2,s;m)$ denote equilibrium-reduced completion and
\[
M_A^*(m)=\max_{(p_1,p_2)\in\cM_A}
\widehat M(p_1,p_2,1;m).
\]
Define
\begin{equation}
\kappa_{\com}^*(m)
=B\sup_{\mu\in\cM_E:Q^{\com}(\mu)>0}
\frac{[\widehat M(\mu)-M_A^*(m)]_+}{Q^{\com}(\mu)}.
\label{eq:kappastar}
\end{equation}

\begin{theorem}[Exact and finite committed-reach no-search threshold]
\label{thm:kappa}
Use the smallest-$s$ policy when outcomes tie.  Then
\[
\mathcal V_E(m;\kappa)=BM_A^*(m)
\quad\Longleftrightarrow\quad
\kappa\ge\kappa_{\com}^*(m).
\]
Moreover, optimal committed outer capacity is weakly decreasing in $\kappa$.
If $\widehat M$ is uniformly Lipschitz in $s$ on the compact policy set,
\[
|\widehat M(p_1,p_2,s)-\widehat M(p_1,p_2,1)|\le L(s-1),
\]
then $\kappa_{\com}^*(m)<\infty$ for every fixed $m>0$.
\end{theorem}

\begin{proof}
For every positive-resource outcome,
\[
B\widehat M-\kappa Q^{\com}\le BM_A^*
\quad\Longleftrightarrow\quad
\kappa\ge B\frac{\widehat M-M_A^*}{Q^{\com}}.
\]
Taking the positive part and supremum proves the threshold, while $s=1$
attains $BM_A^*$.  For resource monotonicity, add the two optimality
inequalities at costs $\kappa_1<\kappa_2$ to obtain
$(\kappa_2-\kappa_1)(Q_1-Q_2)\ge0$.

For finiteness, note $M_A^*(m)>0$.  Whenever $\widehat M>M_A^*$,
completion cannot exceed posting mass, so
$P=1-p_1\ge\widehat M>M_A^*$.  Also $q=\e^{-ma}\ge\e^{-m}$ because
$a\le1$.  The same $(p_1,p_2)$ at $s=1$ is feasible in $A$, so
\[
\widehat M(p_1,p_2,s)-M_A^*
\le\widehat M(p_1,p_2,s)-\widehat M(p_1,p_2,1)
\le L(s-1).
\]
Therefore
\[
\frac{B[\widehat M-M_A^*]}{Q^{\com}}
\le\frac{BLe^m}{mM_A^*(m)}<\infty.
\]
\end{proof}

The uniform Lipschitz condition is a regularity assumption on the
equilibrium-reduced map.  Unique continuity alone is not enough to deliver a
finite ratio bound.

\subsection{Complete-WPBE local search index}

Let $(p_1^A,p_2^A)$ be a unique interior fixed-arrival optimum with a regular
cutoff, $p_2>p_1$, and strict rescue activity at $s=1$.  Evaluate all objects
at this policy and write
\[
g=P\eta_2=\frac{A_2-A_1}{C_2-C_1},
\qquad
K=\frac{p_2}{\beta}-p_1+(1-\omega)\rho_2\e^{-\Lambda_2}.
\]
Let $\mathcal R(a,s)=D_E(a;a)$ and
\[
a_s=-\frac{\mathcal R_s}{\mathcal R_a}
\]
be the cutoff response along the regular WPBE branch.

\begin{theorem}[Local no-search index with equilibrium cutoff response]
\label{thm:local}
At $s=1$, the right derivatives are
\[
M_s=\e^{-ma}m
\left[\rho_2\e^{-\Lambda_2}p_2+K a_s\right],
\]
\[
Q_s^{\com}=P\e^{-ma}m,
\qquad
Q_s^{\ex}=g\e^{-ma}m.
\]
Hence the complete-WPBE local break-even indices are
\[
\boxed{
\kappa_{\mathrm{loc}}^{\com}
=\frac{B}{P}
\left[\rho_2\e^{-\Lambda_2}p_2+K a_s\right] }
\]
and, for executed contacts on a strict rescue branch,
\[
\boxed{
\kappa_{\mathrm{loc}}^{\ex}
=\frac{B}{g}
\left[\rho_2\e^{-\Lambda_2}p_2+K a_s\right] . }
\]
Expansion is locally valuable precisely when the applicable $\kappa$ lies
below its index.  A global conclusion additionally requires a nonincreasing
profile derivative or a single positive-to-negative crossing.
\end{theorem}

\begin{proof}
The implicit-function theorem gives $a_s=-\mathcal R_s/\mathcal R_a$.
Differentiate the rescue branch of \eqref{eq:completion}.  At $s=1$,
$\Lambda_{2,s}=mp_2$ and $\Lambda_{2,a}=-\omega m$.  Collecting the cutoff
terms gives $K$ and the displayed $M_s$.  Differentiating
\eqref{eq:qcom} and \eqref{eq:qex} gives the resource derivatives.  Because
prices are optimized at the fixed-arrival interior optimum, the envelope
theorem removes their first-order response when $s$ is opened.
\end{proof}

The cutoff term is economically important.  Direct reach has two opposing
effects on incumbent waiting.  Greater rescue coverage can move more riders
into rescue, raising waiting value, while more fresh competition lowers the
incumbent's assignment probability.  Thus $a_s$ is not signed globally.

\subsection{Why executed contacts are not the main cost}

Suppose at fixed $(a,p_1,p_2)$ that $p_2>p_1$, $p_2<\beta$,
$C_2(1)>C_1$, and $A_2(1)=A_1$.  Let $x=s-1$.  Then rescue activation gives
\[
M(s)-M(1)=\Theta(x),\qquad
P\eta_2(s)=\Theta(x),\qquad
Q^{\ex}(s)=\Theta(x^2).
\]
If the profiled equilibrium completion retains a strictly positive right
derivative, no finite per-executed-contact cost can implement $s=1$.  This
argument does not apply when $p_1=p_2$, where the rescue mass may already be
first order.  Committed reach instead has
$Q^{\com}(s)=\Theta(x)$ and removes this denominator failure.

\section{Computational falsification and implementation status}

The analytical model is accompanied by an exact Poisson solver that evaluates
every candidate policy at the full type-domain cutoff conditions and then
performs the outer search.  The currently archived numerical tables use the
earlier executed-contact objective, whereas the main theory above recommends
committed reach.  They support the following equilibrium and implementation
claims but are not a committed-cost comparative-static audit:

\begin{itemize}[leftmargin=1.6em]
\item all 20 unit tests pass, including regime identities, independent core
arrivals, incumbent retention, spatial thinning, full best responses, and
outer nestedness;
\item the formal audit contains 99 environments and 297 architecture rows;
all are cutoff-certified and the maximum certified equilibrium-cutoff count
is one;
\item the minimum numerical $E$-minus-$A$ value gap is
$-1.22\times10^{-10}$, which is optimization tolerance, and the maximum is
$0.123759$;
\item none of the 99 archived fixed-arrival optima is close to zero rescue
mass (the minimum is approximately $0.0678$), so the current figures are not
mechanically generated by the executed-contact activation loophole;
\item the maintained-primitives fixed-menu reversal reported above was found
by adversarial search and re-solves the cutoff-WPBE at every thickness.
\end{itemize}

The earlier claim of 48 primitive vectors and 1,344 environments is not
supported by an archived script and result table in the current folder and is
therefore not used in this note.

Two numerical boundary issues remain in the implementation rather than in the
analytical equilibrium engine.  An absolute rider-utility tolerance of
$10^{-12}$ can convert a strict choice into a tie at extremely small $m$, and
rounded rider breakpoints can collapse the conditional support as
$p_1\uparrow1$.  These cases should be replaced by the closed-form upper-
envelope formulas in Lemma~\ref{lem:rider} before the code is used for extreme-
boundary comparative statics.

\section{Interpretation, limitations, and theorem agenda}

The model is internally coherent under its maintained assumptions, but those
assumptions define the scope of every theorem.

\begin{enumerate}[leftmargin=2em]
\item \textbf{Representative-order partial equilibrium.}  Independent local
Poisson pools ignore capacity competition across simultaneous orders.  A
system-wide thick-market claim requires a sparse/mean-field extension.
\item \textbf{Anonymous lottery.}  Search affects the assignment probability
of retained rejectors through competition.  Nearest-driver or history-priority
assignment can change this discipline channel.
\item \textbf{Winner-paid pickup.}  The clean outer thinning relies on the
pickup cost being borne only by the selected driver.  Pre-assignment
relocation requires endogenous entry and congestion.
\item \textbf{Type-independent exit.}  Retention $\omega$ is an availability
shock independent of $c$ and has zero outside payoff.  Selection in exit may
change the posterior and uniqueness proof.
\item \textbf{No rider pickup disutility.}  The rider discounts only across
windows through $\beta$; longer outer pickup distance does not directly lower
rider utility.  Adding that channel would strengthen the case for costly or
interior reach.
\item \textbf{Pass-through price.}  The current objective is not profit.  A
profit paper must separate fare and wage or justify the direct-transfer
institution.
\end{enumerate}

The next genuinely nontrivial theorem is not another endpoint argument.  It is
a primitive region in which the fully profiled scope derivative single-crosses
or in which the separation margins can be verified analytically.  Such a
result could deliver a scarce-intermediate-thick search window after full WPBE
response and reoptimization.  Until then, the correct hierarchy is:

\begin{enumerate}[leftmargin=2em]
\item unique incumbent cutoff and exact completion aggregation: proved;
\item endpoint collapse and existence of a positive interior maximizer:
proved;
\item finite committed-reach no-search threshold under uniform regularity:
proved;
\item ordered peaks with frozen hazards: proved as a benchmark;
\item ordered peaks after full WPBE reoptimization: numerical conjecture,
not a theorem;
\item global unique peaks and monotone $s^*(m)$: false or unproved.
\end{enumerate}

\section{Conclusion}

Universal rejection is the organizing economic event.  It screens the
incumbent pool, informs the rider, and activates a recovery market in which
one common transaction price changes both retained-driver eligibility and
fresh-driver willingness.  Fixed-footprint arrivals and geographic expansion
restore supply through different channels: the former changes the supply
environment in every terminal action, while the latter changes who first sees
rescue and adds pickup friction.  The platform's menu must therefore be solved
as an equilibrium-constrained posted mechanism.

The strongest currently valid theoretical result is a chain rather than a
single closed form: conditional Poisson supply, unique cutoff-WPBE outcome,
exact rider aggregation, full-design endpoint collapse, and a committed-reach
cost threshold.  The elegant closed-form peak order identifies an important
Poisson force but is not universal once equilibrium composition and policy
reoptimization are allowed.  That boundary is a result of the adversarial
audit, not a defect to conceal.

\appendix

\section{Proof of the unique-cutoff theorem}
\label{app:unique}

For uniform costs and $a\le p_1\le p_j$,
\[
\Lambda_j(a)=\omega m(p_j-a)+e_j,
\qquad \Lambda_{j,a}=-\omega m.
\]
In all three architectures, $e_2\ge e_1$, so
\[
\Lambda_2(a)-\Lambda_1(a)
=\omega m(p_2-p_1)+(e_2-e_1)\ge0
\]
is independent of $a$.

\paragraph{Step 1: type single crossing and existence.}
Away from the kinks $p_1,p_2$,
\[
\frac{\partial D_r(c;a)}{\partial c}
=-h(ma)+\delta\omega\e^{-ma}
\sum_{j:c<p_j}\eta_jh(\Lambda_j)
\le-h(ma)+\e^{-ma}\le0.
\]
The final inequality uses $h(x)\ge\e^{-x}$, equivalent to
$\e^x-1\ge x$.  Thus the type payoff difference is globally single crossing.
Let $R(a)=D_r(a;a)$.  At $a=p_1$, immediate surplus is zero and waiting
surplus is nonnegative, so $R(p_1)\le0$.  If $R(0)\le0$, $a=0$ is a cutoff
best response.  Otherwise continuity and the intermediate value theorem give
a root in $(0,p_1]$, and type single crossing supplies the full best-response
inequalities.

\paragraph{Step 2: derivative bounds.}
On the active domain $p_1<\beta$, total rider continuation mass is
\[
\ell=\frac{1-p_1/\beta}{1-p_1}\le1.
\]
On a strict repeat-rescue branch,
\[
\eta_{2,a}=
\frac{\omega m(p_2-p_1)}
{\beta(1-p_1)(C_2-C_1)}\ge0,
\qquad \eta_1=\ell-\eta_2.
\]
Define
\[
\mathcal A(a)=h(ma)(p_1-a),
\qquad
Y_j(a)=\e^{-ma}h(\Lambda_j(a))(p_j-a).
\]
The integral representation $h(z)=\int_0^1\e^{-zt}dt$ gives
$h'<0$, $h+h'>0$, and $h(z)+zh'(z)=\e^{-z}$.  Put
$y_j=\omega m(p_j-a)\le\Lambda_j$.  Then
\begin{align*}
-\omega\e^{ma}Y_j'
&=\omega h(\Lambda_j)
+y_j[h(\Lambda_j)+\omega h'(\Lambda_j)]\\
&\le \omega h(\Lambda_j)
+\Lambda_j[h(\Lambda_j)+\omega h'(\Lambda_j)]\\
&=1-(1-\omega)\e^{-\Lambda_j}\le1.
\end{align*}
Hence
\begin{equation}
-\omega Y_j'(a)\le\e^{-ma}.
\label{eq:Ybound}
\end{equation}
Also
\begin{equation}
-\mathcal A'(a)
=h(ma)-m(p_1-a)h'(ma)>\e^{-ma}.
\label{eq:Abound}
\end{equation}

\paragraph{Step 3: every interior root is a strict downward crossing.}
At any nondegenerate interior residual root on an active rescue branch,
\[
h(\Lambda_2)(p_2-a)>h(\Lambda_1)(p_1-a).
\]
Otherwise discounted waiting is at most
$\delta\omega\e^{-ma}\ell(p_1-a)$, strictly below
$h(ma)(p_1-a)$, contradicting indifference.  Therefore $Y_2>Y_1$ at a root.
Let
\[
\mathcal T(a)=(\ell-\eta_2)Y_1+\eta_2Y_2.
\]
The composition derivative $\eta_{2,a}(Y_2-Y_1)$ is nonnegative.  Applying
\eqref{eq:Ybound} gives
\[
-\omega\mathcal T'(a)\le\ell\e^{-ma}.
\]
Since $R=\mathcal A-\delta\omega\mathcal T$, equation
\eqref{eq:Abound} yields at every interior root
\[
R'(a)<-\e^{-ma}+\delta\ell\e^{-ma}\le0.
\]

Before rescue activates, only repeat at $p_1$ is relevant and, for every
$a<p_1$,
\[
R(a)=(p_1-a)
\left[h(ma)-\delta\omega\ell\e^{-ma}h(\Lambda_1)\right]>0.
\]
Thus the rescue-activation point cannot itself be a root.  Starting from the
positive inactive branch, a continuous scalar residual whose every active
root crosses strictly downward can have at most one root.  Existence therefore
implies uniqueness.

The boundary cases follow directly: if $p_1=0$, the cutoff domain is the
singleton $a=0$; if $p_1\ge\beta$ or $\omega=0$, waiting has no positive
value and $a=p_1$; if $p_2=p_1$, delay and weakly worse assignment dominate
waiting, again giving $a=p_1$.  The point $p_1=1$ is the continuous zero-
demand closure.

For continuity, take any convergent sequence of policies and primitives and
their unique cutoffs.  Every subsequential cutoff limit satisfies the closed
best-response inequalities of the limiting problem.  Uniqueness identifies
that limit, so the full sequence converges.  Lemma~\ref{lem:rider} then gives
continuity of completion.  Under the deterministic closure,
$Q^{\com}$ is continuous as well.

\section{Audited claim ledger}

\begin{longtable}{@{}p{0.48\textwidth}p{0.46\textwidth}@{}}
\toprule
Claim & Status after adversarial audit \\
\midrule
\endfirsthead
\toprule
Claim & Status after adversarial audit \\
\midrule
\endhead
Conditional Poisson supply after universal rejection & Proved under marked
Poisson costs and type-independent retention. \\
Unique full WPBE in every action label & Too strong.  Unique incumbent cutoff
and completion; payoff-equivalent labels may differ. \\
Exact rider completion formula & Proved globally with
$\max\{0,A_1,A_2\}$. \\
Outer optimum exists & Proved for compact policies and the continuous
committed-reach objective. \\
Both incremental gaps vanish in thin and thick limits & Proved after full
WPBE response and reoptimization. \\
Any positive gap has an interior maximum & Proved; uniqueness is not implied. \\
Expanded-search value always peaks first & False as a universal fixed-menu
claim; maintained-primitives numerical reversal exists. \\
Frozen-hazard Poisson peaks are ordered & Proved with closed forms. \\
Uniform separation preserves the order & Proved as a robustness certificate,
not a primitive condition. \\
Finite no-search threshold under committed reach & Proved under uniform
Lipschitz regularity. \\
Finite no-search threshold under executed contacts & False at a qualifying
rescue-activation boundary. \\
$s^*(m)$ is globally monotone & False in numerical examples; an off-on-off
pattern can arise. \\
The current objective is profit or welfare & False.  It is completion value
net of a search-resource cost. \\
$I\subset A\subset E$ as mechanism sets & False.  Only $A\subset E$ through
$s=1$; $I$ and $A$ are supply environments. \\
\bottomrule
\end{longtable}

\begin{thebibliography}{99}

\bibitem{HuHuZhu2022}
Hu, B., Hu, M., and Zhu, H. (2022).
Surge pricing and two-sided temporal responses in ride hailing.
\emph{Manufacturing \& Service Operations Management}, 24(1), 91--109.
\href{https://doi.org/10.1287/msom.2020.0960}{doi:10.1287/msom.2020.0960}.

\bibitem{QinYangLiu2025}
Qin, X., Yang, H., and Liu, Y. (2025).
A two-round broadcasting matching mechanism in ride-sourcing markets:
Implication and optimization.
\emph{Transportation Research Part C}, 180, 105318.
\href{https://doi.org/10.1016/j.trc.2025.105318}{doi:10.1016/j.trc.2025.105318}.

\bibitem{LyftI2026}
Ekbatani, F., Niazadeh, R., Golari, M., et al. (2026a).
Non-exclusive notifications for ride-hailing at Lyft I: Single-cycle
approximation algorithms.
\href{https://arxiv.org/abs/2603.21533}{arXiv:2603.21533}.

\bibitem{LyftII2026}
Ekbatani, F., Niazadeh, R., Golari, M., et al. (2026b).
Non-exclusive notifications for ride-hailing at Lyft II: Simulations and
marketplace analysis.
\href{https://arxiv.org/abs/2603.21531}{arXiv:2603.21531}.

\bibitem{WangZhangFengKe2026}
Wang, C., Zhang, K., Feng, S., and Ke, J. (2026).
Pricing and matching for ride-hailing markets under the broadcasting
mechanism.
SSRN Working Paper 7240784.
\href{https://doi.org/10.2139/ssrn.7240784}{doi:10.2139/ssrn.7240784}.

\bibitem{XuYan2026}
Xu, A., and Yan, C. (2026).
Dispatching and pricing in two-sided spatial queues.
\emph{Transportation Science}, Articles in Advance.
\href{https://doi.org/10.1287/trsc.2025.0319}{doi:10.1287/trsc.2025.0319}.

\bibitem{ChenHu2020}
Chen, Y., and Hu, M. (2020).
Pricing and matching with forward-looking buyers and sellers.
\emph{Manufacturing \& Service Operations Management}, 22(4), 717--734.
\href{https://doi.org/10.1287/msom.2018.0769}{doi:10.1287/msom.2018.0769}.

\bibitem{CastroEtAl2022}
Castro, F., Ma, H., Nazerzadeh, H., and Yan, C. (2022).
Randomized FIFO mechanisms.
\emph{Proceedings of the 23rd ACM Conference on Economics and Computation}.
\href{https://doi.org/10.1145/3490486.3538353}{doi:10.1145/3490486.3538353}.

\bibitem{GargNazerzadeh2022}
Garg, N., and Nazerzadeh, H. (2022).
Driver surge pricing.
\emph{Management Science}, 68(5), 3219--3235.
\href{https://doi.org/10.1287/mnsc.2021.4058}{doi:10.1287/mnsc.2021.4058}.

\bibitem{YangQinKeYe2020}
Yang, H., Qin, X., Ke, J., and Ye, J. (2020).
Optimizing matching time interval and matching radius in on-demand
ride-sourcing markets.
\emph{Transportation Research Part B}, 131, 84--105.
\href{https://doi.org/10.1016/j.trb.2019.11.005}{doi:10.1016/j.trb.2019.11.005}.

\bibitem{VarmaEtAl2023}
Varma, S. M., Bumpensanti, P., Maguluri, S. T., and Wang, H. (2023).
Dynamic pricing and matching for two-sided queues.
\emph{Operations Research}, 71(1), 83--100.
\href{https://doi.org/10.1287/opre.2021.2233}{doi:10.1287/opre.2021.2233}.

\bibitem{AfecheEtAl2023}
Af\`eche, P., Liu, Z., and Maglaras, C. (2023).
Ride-hailing networks with strategic drivers: The impact of platform control
capabilities on performance.
\emph{Manufacturing \& Service Operations Management}.
\href{https://doi.org/10.1287/msom.2023.1221}{doi:10.1287/msom.2023.1221}.

\bibitem{TripathyEtAl2023}
Tripathy, M., Bai, J., and Heese, H. S. (2023).
Driver collusion in ride-hailing platforms.
\emph{Decision Sciences}, 54(4), 434--446.
\href{https://doi.org/10.1111/deci.12561}{doi:10.1111/deci.12561}.

\bibitem{TuEtAl2026}
Tu, W., et al. (2026).
Effect of trip attributes on ridehailing driver trip request acceptance.
\emph{International Journal of Sustainable Transportation}.
\href{https://doi.org/10.1080/15568318.2025.2580366}{doi:10.1080/15568318.2025.2580366}.

\end{thebibliography}

\end{document}
